%%FILE blueprint/src/macros/common.tex
% In this file you should put all LaTeX macros and settings to be used both by
% the pdf version and the web version.
% This should be most of your macros.

% The theorem-like environments defined below are those that appear by default
% in the dependency graph. See the README of leanblueprint if you need help to
% customize this.
% The configuration below use the theorem counter for all those environments
% (this is what the [theorem] arguments mean) and never resets it.
% If you want for instance to number them within chapters then you can add
% [chapter] at the end of the next line.
\newtheorem{theorem}{Theorem}[chapter] % Restart numbering in each chapter
\newtheorem{proposition}[theorem]{Proposition}
\newtheorem{lemma}[theorem]{Lemma}
\newtheorem{corollary}[theorem]{Corollary}

\theoremstyle{definition}
\newtheorem{definition}[theorem]{Definition}

\DeclareMathOperator{\Od}{Od}
\DeclareMathOperator{\conv}{conv}
\DeclareMathOperator{\per}{per}
% \cap is already taken, so we name the command `\mycap` instead
\DeclareMathOperator{\mycap}{cap}
\DeclareMathOperator{\supp}{supp}
\DeclareMathOperator{\Prop}{Prop}
% \cr is already taken, so we name the command `\mycr` instead
\DeclareMathOperator{\mycr}{cr}
\DeclareMathOperator{\ord}{ord}
\DeclareMathOperator{\trace}{trace}
\DeclareMathOperator{\sign}{sign}

%%FILE blueprint/src/macros/web.tex
% In this file you should put macros to be used only by
% the web version. Of course they should have a corresponding
% version in macros/print.tex.
% Typically the printed version could have more fancy decorations.
% This will probably be a very short file.
%%FILE blueprint/src/macros/print.tex
% In this file you should put macros to be used only by
% the printed version. Of course they should have a corresponding
% version in macros/web.tex.
% Typically the printed version could have more fancy decorations.
% This should be a very short file.
%
% This file starts with dummy macros that ensure the pdf
% compiler will ignore macros provided by plasTeX that make
% sense only for the web version, such as dependency graph
% macros.


% Dummy macros that make sense only for web version.
\newcommand{\lean}[1]{}
\newcommand{\discussion}[1]{}
\newcommand{\leanok}{}
\newcommand{\mathlibok}{}
\newcommand{\notready}{}
% Make sure that arguments of \uses and \proves are real labels, by using invisible refs:
% latex prints a warning if the label is not defined, but nothing is shown in the pdf file.
% It uses LaTeX3 programming, this is why we use the expl3 package.
\ExplSyntaxOn
\NewDocumentCommand{\uses}{m}
 {\clist_map_inline:nn{#1}{\vphantom{\ref{##1}}}%
  \ignorespaces}
\NewDocumentCommand{\proves}{m}
 {\clist_map_inline:nn{#1}{\vphantom{\ref{##1}}}%
  \ignorespaces}
\ExplSyntaxOff

% Redefine \mathbb command as in amsfonts
\AtBeginDocument{%
  \let\mathbb\relax
  \DeclareMathAlphabet{\mathbb}{U}{msb}{m}{n}%
}

% Read authors.md file, puts \and between author names.
\newread\authorsin
\newtoks\authortoks
\newif\iffirstauthor

\def\readauthors#1{%
  \authortoks={}%
  \firstauthortrue
  \openin\authorsin=#1\relax
  \loop
    \ifeof\authorsin
      \let\next=\relax
    \else
      \read\authorsin to \thisauthor
      \ifeof\authorsin
        \let\next=\relax
      \else
        \iffirstauthor
          \edef\temp{\thisauthor}%
          \firstauthorfalse
        \else
          \edef\temp{\the\authortoks\space\noexpand\and\space\thisauthor}%
        \fi
        \authortoks=\expandafter{\temp}%
        \let\next=\repeat
      \fi
    \fi
  \next
  \closein\authorsin
  \the\authortoks
}